\section{Morphisms of ringed spaces and modules}
\label{section-ringed-spaces-functoriality-modules}

\noindent
We have now introduced enough notation so that we are able to
define the pullback and pushforward of modules along a morphism
of ringed spaces.

\begin{definition}
\label{definition-pushforward}
Let $(f, f^\sharp) : (X, \mathcal{O}_X) \to (Y, \mathcal{O}_Y)$
be a morphism of ringed spaces.
\begin{enumerate}
\item Let $\mathcal{F}$ be a sheaf of $\mathcal{O}_X$-modules.
We define the {\it pushforward} of $\mathcal{F}$ as the
sheaf of $\mathcal{O}_Y$-modules which as a sheaf
of abelian groups equals $f_*\mathcal{F}$ and with
module structure given by the restriction
via $f^\sharp : \mathcal{O}_Y \to f_*\mathcal{O}_X$
of the module structure given
in Lemma \ref{lemma-pushforward-module}.
\item Let $\mathcal{G}$ be a sheaf of $\mathcal{O}_Y$-modules.
We define the {\it pullback} $f^*\mathcal{G}$ to be the
sheaf of $\mathcal{O}_X$-modules defined by the formula
$$
f^*\mathcal{G}
=
\mathcal{O}_X \otimes_{f^{-1}\mathcal{O}_Y} f^{-1}\mathcal{G}
$$
where the ring map $f^{-1}\mathcal{O}_Y \to \mathcal{O}_X$
is the map corresponding to $f^\sharp$, and where the  module
structure is given by Lemma \ref{lemma-pullback-module}.
\end{enumerate}
\end{definition}

\noindent
Thus we have defined functors
\begin{eqnarray*}
f_* : \textit{Mod}(\mathcal{O}_X)
& \longrightarrow &
\textit{Mod}(\mathcal{O}_Y) \\
f^* : \textit{Mod}(\mathcal{O}_Y)
& \longrightarrow &
\textit{Mod}(\mathcal{O}_X)
\end{eqnarray*}
The final result on these functors is that they are indeed
adjoint as expected.

\begin{lemma}
\label{lemma-adjoint-pullback-pushforward-modules}
Let $(f, f^\sharp) : (X, \mathcal{O}_X) \to (Y, \mathcal{O}_Y)$
be a morphism of ringed spaces.
Let $\mathcal{F}$ be a sheaf of $\mathcal{O}_X$-modules.
Let $\mathcal{G}$ be a sheaf of $\mathcal{O}_Y$-modules.
There is a canonical bijection
$$
\Hom_{\mathcal{O}_X}(f^*\mathcal{G}, \mathcal{F})
=
\Hom_{\mathcal{O}_Y}(\mathcal{G}, f_*\mathcal{F}).
$$
In other words: the functor $f^*$ is the left adjoint to
$f_*$.
\end{lemma}

\begin{proof}
This follows from the work we did before:
\begin{eqnarray*}
\Hom_{\mathcal{O}_X}(f^*\mathcal{G}, \mathcal{F})
& = &
\Mor_{\textit{Mod}(\mathcal{O}_X)}(
\mathcal{O}_X \otimes_{f^{-1}\mathcal{O}_Y} f^{-1}\mathcal{G}, \mathcal{F}) \\
& = &
\Mor_{\textit{Mod}(f^{-1}\mathcal{O}_Y)}(
f^{-1}\mathcal{G}, \mathcal{F}_{f^{-1}\mathcal{O}_Y}) \\
& = &
\Hom_{\mathcal{O}_Y}(\mathcal{G}, f_*\mathcal{F}).
\end{eqnarray*}
Here we use Lemmas \ref{lemma-adjointness-tensor-restrict}
and \ref{lemma-adjoint-push-pull-modules}.
\end{proof}

\begin{lemma}
\label{lemma-push-pull-composition-modules}
Let $f : X \to Y$ and $g : Y \to Z$ be morphisms of ringed spaces.
The functors $(g \circ f)_*$ and $g_* \circ f_*$ are equal.
There is a canonical isomorphism of functors
$(g \circ f)^* \cong f^* \circ g^*$.
\end{lemma}

\begin{proof}
The result on pushforwards is a consequence of Lemma
\ref{lemma-pushforward-composition} and our definitions.
The result on pullbacks follows from this by the same
argument as in the proof of Lemma \ref{lemma-pullback-composition}.
\end{proof}


\noindent
Given a morphism of ringed spaces
$(f, f^\sharp) : (X, \mathcal{O}_X) \to (Y, \mathcal{O}_Y)$,
and a sheaf of $\mathcal{O}_X$-modules $\mathcal{F}$,
a sheaf of $\mathcal{O}_Y$-modules $\mathcal{G}$ on $Y$,
the notion of an {\it $f$-map $\varphi : \mathcal{G} \to \mathcal{F}$
of sheaves of modules} makes sense. We can just define
it as an $f$-map $\varphi : \mathcal{G} \to \mathcal{F}$
of abelian sheaves (see Definition \ref{definition-f-map} and
Lemma \ref{lemma-f-map}) such that for all open $V \subset Y$ the map
$$
\mathcal{G}(V) \longrightarrow \mathcal{F}(f^{-1}V)
$$
is an $\mathcal{O}_Y(V)$-module map. Here we think of
$\mathcal{F}(f^{-1}V)$ as an $\mathcal{O}_Y(V)$-module
via the map $f^\sharp_V : \mathcal{O}_Y(V) \to \mathcal{O}_X(f^{-1}V)$.
The set of $f$-maps between
$\mathcal{G}$ and $\mathcal{F}$ will be in canonical bijection
with the sets
$\Mor_{\textit{Mod}(\mathcal{O}_X)}(f^*\mathcal{G}, \mathcal{F})$
and
$\Mor_{\textit{Mod}(\mathcal{O}_Y)}(\mathcal{G}, f_*\mathcal{F})$.
See above.

\medskip\noindent
Composition of $f$-maps is defined in exactly the
same manner as in the case of $f$-maps of sheaves of
sets. In addition, given an $f$-map $\mathcal{G} \to \mathcal{F}$
as above, and $x \in X$ the induced map on stalks
$$
\varphi_x : \mathcal{G}_{f(x)} \longrightarrow \mathcal{F}_x
$$
is an $\mathcal{O}_{Y, f(x)}$-module map where the
$\mathcal{O}_{Y, f(x)}$-module structure on $\mathcal{F}_x$
comes from the $\mathcal{O}_{X, x}$-module structure via
the map $f^\sharp_x : \mathcal{O}_{Y, f(x)} \to \mathcal{O}_{X, x}$.
Here is a related lemma.

\begin{lemma}
\label{lemma-stalk-pullback-modules}
Let $(f, f^\sharp) : (X, \mathcal{O}_X) \to (Y, \mathcal{O}_Y)$
be a morphism of ringed spaces.
Let $\mathcal{G}$ be a sheaf of $\mathcal{O}_Y$-modules.
Let $x \in X$. Then
$$
(f^*\mathcal{G})_x =
\mathcal{G}_{f(x)}
\otimes_{\mathcal{O}_{Y, f(x)}}
\mathcal{O}_{X, x}
$$
as $\mathcal{O}_{X, x}$-modules where the tensor product on the right
uses $f^\sharp_x : \mathcal{O}_{Y, f(x)} \to \mathcal{O}_{X, x}$.
\end{lemma}

\begin{proof}
This follows from Lemma \ref{lemma-stalk-tensor-sheaf-modules}
and the identification of the stalks of pullback sheaves
at $x$ with the corresponding stalks at $f(x)$. See the
formulae in Section \ref{section-presheaves-structures-functorial}
for example.
\end{proof}



